\documentclass[a4paper,10pt,fleqn]{article}
\usepackage[margin=1in]{geometry}
\usepackage{amssymb}
\usepackage{adjustbox}
\usepackage{natbib}
\usepackage[colorlinks=true,linkcolor=blue,citecolor=blue,urlcolor=blue]{hyperref}
\usepackage{fuzz}
\usepackage{schemapk}
\usepackage{zed-maths}
\usepackage{zed-proof}
\newdimen\savedleftskip
\begin{document}

\section*{String Literals in Z}
\addcontentsline{toc}{section}{String Literals in Z}

\noindent In database modelling, constants often have string values. The txt2tex lexer recognises single-quoted strings as literal values distinct from identifiers.

\bigskip

\noindent A string literal like 'sunk' or 'survived' is a concrete value. When used in Z expressions, it renders with Z-convention quoting: backtick open, apostrophe close.

\bigskip

\noindent In this TEXT block, the apostrophes appear as written. In a Z expression block such as an axdef or schema predicate, 'sunk' would render as the Z string form.

\bigskip

\end{document}